% ID: SEL-D0044-OB-E02
% ESTADO: validada. S045, cotejada con examen y criterios.
\textbf{a)} Cada tramo es continuo y derivable en su interior.
En $1$, el valor y ambos límites valen 4. Por tanto es continua
en todo $\mathbb R$. Las pendientes laterales en $1$ son
$2(1+1)=4$ y $-4/1^2=-4$. Como difieren, es derivable
en $\mathbb R\setminus\{1\}$.

\textbf{b)} Al tender $x$ a $+\infty$, $4/x\to0$.
La asíntota horizontal es $\boxed{y=0}$. El denominador no se
anula en $x>1$ y la función es continua en todo su dominio.
No hay asíntotas verticales. Hacia $-\infty$ domina el
término cuadrático, así que allí no hay asíntota recta.

\textbf{c)} $f(2)=2$ y $f'(2)=-1$. La tangente
$y-2=-(x-2)$ queda $\boxed{y=4-x}$.
